\section{Constitutive evaluation and dimensional realization}
\label{iii:sec:evaluacion}

Evaluation now brings together the preceding constructions in their effective order.
The regional expansions of $\pi$, $\varphi$, and $e$ arise from nonadic
refinement; the register $K$ preserves its twelve positions. Its periodic readout
determines the positional section of $\alpha$. The action character
and regional return act on it, followed by the angular pair and, finally, the
constitutive response. This evaluation follows the selection of
rules: none of its digits is used to choose a channel.

\subsection{Evaluation of the internal chain}

To reproduce the table, one uses the thousand-decimal regional expansions
generated by the nonadic certificate and the integer $N_K$ in
\eqref{eq:k-numerador}. The readout used is
\[
 \alpha=\pi+e-\varphi-4-\frac{N_K}{1000^{12}-1}.
\]
The finite twelve-block section and the prolonged section have their
definitions and exact difference in the register chapter. Here the
latter is held fixed throughout evaluation. Readouts are not interchanged
as precision increases, nor is the analytic operator
of the second route replaced by a decimal coincidence with the first.

\begin{longtable}{@{}ll@{}}
\toprule
Object & Decimal evaluation (24 significant digits)\\
\midrule\endhead
$\alpha$ & $0.00729735256928380099728511\ldots$ \\
$H_5$ & $1.05457181764615446962582\ldots$ \\
$\eta_{\rm ret}$ & $1.05457181691503906113369\ldots$ \\
$A_{\deg}$ & $7.29735256928380099728511\ldots$ \\
$C^*_{\deg}$ & $2.12373833896864572426195\ldots$ \\
$q_+$ & $0.848377942576404374945737\ldots$ \\
$q_-$ & $0.913660151150557285295519\ldots$ \\
$r_+$ & $0.999999628548249363178695\ldots$ \\
$r_-$ & $0.999724137189518364131517\ldots$ \\
$\widehat\varepsilon$ & $0.999999257096636702760442\ldots$ \\
$\widehat\mu$ & $0.999448350479326935089982\ldots$ \\
$\widehat Z$ & $0.999724508538937215455554\ldots$ \\
$\widehat c$ & $1.00027631048615752610639\ldots$ \\
$\gamma$ & $0.274007672694459274747255\ldots$ \\
$G_{\rm Cat}$ & $0.915965594177219015054604\ldots$ \\
$\rho$ & $0.267949192431122706472554\ldots$ \\
\bottomrule
\end{longtable}


The calculation uses 120 working digits; 24 significant
digits are displayed to facilitate comparison between quantities of different
orders. The 120 digits are an evaluation precision, not an experimental
uncertainty or, by themselves, an interval error certificate.
The package preserves the generated inputs, their hashes, the operations, and
the residuals of the checked identities. The algebraic proofs in the
main text establish these identities independently of this table.

The meaning of the final four constitutive coordinates is preserved
under evaluation: $\widehat\varepsilon$ and $\widehat\mu$ are the quadratic
responses per sheet, $\widehat Z$ is their relative orientation, and
$\widehat c$ is the inverse of their product. The proximity of some digits to
unity expresses the small defect of the $3$--$4$ quotient on these
channels. Its origin is seen exactly in
\[
 1-f(s)=\frac{s^3}{1+s+s^2+s^3}.
\]
The formula allows the response to be studied before a dimensional
quantity is assigned to it: the sign of the defect, its cubic order, and its monotonicity
are consequences of the same operator.

\subsection{Unit charts and covariance}

Let $\mathcal L_S$, $\mathcal L_Q$, and $\mathcal L_C$ be the action,
charge, and velocity lines. A positive physical chart evaluates the sections
$u_S$, $u_Q$, and $u_C$ in action, charge, and velocity units. The realization
\eqref{iii:eq:dimensions} then produces an impedance, a velocity,
a permeability, and a permittivity, with their respective dimensions.
The coordinates in the preceding table remain dimensionless.

\begin{proposition}[Covariance of the constitutive realization]
If the reference sections change by
$u_S'=a u_S$, $u_Q'=b u_Q$, $u_C'=d u_C$, with $a,b,d>0$, the
coordinates representing the same quantities satisfy
\[
 \widehat Z'=a^{-1}b^2\widehat Z,\qquad
 \widehat c'=d^{-1}\widehat c,\qquad
 \widehat\mu'=a^{-1}db^2\widehat\mu,\qquad
 \widehat\varepsilon'=adb^{-2}\widehat\varepsilon.
\]
The relations $\widehat\mu'\widehat\varepsilon'=(\widehat c')^{-2}$
and $\widehat\mu'/\widehat\varepsilon'=(\widehat Z')^2$ are preserved.
\end{proposition}
\begin{proof}
Equate each quantity in \eqref{iii:eq:dimensions} to its expression
in the new basis and solve for its coordinate. The factors in the product
are $d^2$, and those in the quotient are $a^{-2}b^4$, exactly those required
by the transformed coordinates of $c$ and $Z$.
\end{proof}

This covariance distinguishes two acts: constructing a section and expressing
that section in chosen units. The spectral theorem fixes the coordinates
in the declared internal normalization. An SI comparison must also exhibit
the evaluation of the three reference sections. Fitting them
using the quantities to be reproduced would calibrate
the chart, not independently verify those quantities.

To reconstruct the channels from coordinates expressed in the primed
chart, undo that change before using the cubic root:
\begin{equation}
 \widehat Z=ab^{-2}\widehat Z',\quad
 \widehat c=d\widehat c',\quad
 \widehat\mu=ad^{-1}b^{-2}\widehat\mu',\quad
 \widehat\varepsilon=a^{-1}d^{-1}b^2\widehat\varepsilon'.
 \label{iii:eq:undo-units}
\end{equation}
For example, the inputs to the spectral inverse are
$r_+=(\widehat Z\widehat c)^{-1/2}$ and
$r_-=(\widehat Z/\widehat c)^{1/2}$, calculated after
\eqref{iii:eq:undo-units}. Directly inserting the primed coordinates
into those expressions would change the operator normalization and could
take them outside the interval $(3/4,1)$. Dimensional covariance preserves
physical quantities; cubic inversion reconstructs the transport
in the internal chart in which it was defined.

\subsection{Identity checks and scope of comparison}

Reproducible comparison is organized by relations. First,
$\widehat\mu\widehat\varepsilon\widehat c^2=1$ and
$\widehat\mu/(\widehat\varepsilon\widehat Z^2)=1$ check
constitutive consistency. Second, cubic inversion recovers
the original arguments and angular coordinates. Finally, evaluating
Barbero through the angular functional and through functional calculus
of the response agrees by the identity proved in the previous chapter.
Interchanging sheets must reverse the oriented contribution and preserve
the common mode; this check detects a loss of labeling that scalar
products would not reveal.

Comparison of the reduced Planck constant with reference
tables is preserved in the preceding action development. For the
four vacuum quantities, the present table evaluates their
internal coordinates. It is not presented as a table of SI matches:
that additional assertion requires an independent numerical chart for
$u_S,u_Q,u_C$. The proved constitutive result, including its inversion,
remains complete on its domain and explicitly preserves the dimensional
interface required to make that comparison.
